\chapter{Introduction}\label{ch:intro}

This book is a user manual for Dylan, the object-oriented dynamic language
that Apple Computer described in 1992, and for the two implementations of
it that this project provides: an interactive implementation written in
Common Lisp and a compiler that translates Dylan programs into C. It
teaches the language from the beginning, explains how to use each
implementation, and documents where the implementations depart from the
1992 definition.

This chapter describes the language in broad terms, introduces the two
implementations and when to choose each, gives a short history of the
project, and explains how the rest of the book is organized and what its
typographic conventions mean.

\section{What Dylan is}\label{sec:intro-what}

\index{Dylan!1992 definition}
Dylan was designed at Apple Computer's Eastern Research and Technology
laboratory and published in April 1992 as \emph{Dylan: an object-oriented
dynamic language}. That book, which this manual calls the \term{DIRM}, is
the standard both implementations follow. The language it defines is the
original Dylan, written in a fully parenthesized \term{prefix syntax} like
that of Scheme and Common Lisp. Later versions of Dylan moved to an infix
syntax that looks more like Pascal or C; that later language is not the
subject of this book. Everything here is the 1992 prefix language.

A short piece of Dylan shows its flavor:

\begin{dylan}
(define-class <point> (<object>)
  (x init-value: 0 init-keyword: x:)
  (y init-value: 0 init-keyword: y:))

(define-method norm ((p <point>))
  (sqrt (+ (* (x p) (x p)) (* (y p) (y p)))))

(norm (make <point> x: 3 y: 4))   ; => 5.0
\end{dylan}

The language's main ideas are few, and they fit together closely.

\paragraph{Objects all the way down.} Every value in Dylan is an object,
and every object is an instance of a class. Integers, characters,
strings, lists, functions and classes themselves are all objects, with
classes such as \code{<integer>}, \code{<string>} and \code{<class>}.
There are no ``primitive'' values that stand outside the object system,
so any value can be passed to any function, stored in any collection, and
used to select a method.

\paragraph{Generic functions and multiple dispatch.}\index{generic function}
In most object-oriented languages a method belongs to a class and is
chosen by the class of a single receiver object. In Dylan, methods belong
to \term{generic functions}. A generic function is an ordinary function
that holds a set of methods; when it is called, it looks at the classes of
\emph{all} of its required arguments and runs the most specific method that
applies. This is called \term{multiple dispatch}. A method may also be
specialized on one particular object (a \term{singleton}) rather than on a
class.

\paragraph{Classes with multiple inheritance.} A class names one or more
superclasses and describes \term{slots}, the named fields each instance
holds. Dylan defines the getter and setter functions for each slot
automatically, and those functions are themselves generic functions, so
slot access and method calls look exactly alike. A class may inherit from
several superclasses at once.

\paragraph{First-class functions.} Functions are objects. They can be
created anonymously with \code{method}, passed as arguments, returned as
results, and stored in data structures. A function created inside another
function remembers the variables around it (it is a \term{closure}). The
library includes a set of functional operations, such as \code{curry},
\code{compose} and \code{complement}, that build new functions from old
ones.

\paragraph{Conditions.} Errors and other exceptional situations are
represented by \term{condition} objects, which are instances of classes
such as \code{<error>} and \code{<warning>}. A program signals a condition
and a dynamically established handler decides what to do with it: return
a value, take a non-local exit, or decline and let an outer handler try.
Because conditions are objects, handlers select them by class, just as
methods do.

\paragraph{Collections.} Dylan has a rich, uniform library of collections:
lists, vectors, strings, stretchy vectors, deques, ranges, arrays and
tables. All of them share one iteration protocol, so the same operations
(\code{map}, \code{reduce}, \code{do}, \code{choose}, \code{member?},
\code{for-each} and the rest) work on every kind of collection, and new
collection classes can join in by implementing the protocol.

\paragraph{A dynamic language that can be compiled well.} Dylan programs
need no type declarations, and in the 1992 definition classes and methods
can be added and changed while a program runs. The language also offers
ways to say that parts of a program will not change, such as sealing a
class, so that a compiler can produce fast code. The two implementations
described here sit at the two ends of this range.

\section{The two implementations}\label{sec:intro-implementations}

\index{implementations!choosing}
This project provides two separate implementations of the same language.
They share the reader (the code that turns program text into data), and
they are tested against the same examples from the DIRM, but otherwise
they are independent programs.

\paragraph{The interactive implementation.}\index{interactive implementation}
The first is written in Common Lisp and runs inside SBCL. Its source is in
\file{src/}, and this book calls it \emph{the interactive
implementation} or \emph{the Common Lisp implementation}. It is not an
interpreter: it compiles each Dylan form into Common Lisp as the form is
entered and lets the Lisp compiler run it. You use it in three ways: at an
interactive read-eval-print loop, the command \cmd{dylan}; by
compiling a file of Dylan into a file of Common Lisp with
\code{dylan:compile-file}; and by loading such compiled files into a Lisp
that runs them. It implements the whole language of the DIRM, including
the parts that let a program change its own classes and generic functions
as it runs (\code{add-method}, \code{remove-method}, \code{add-slot} and
their relatives), and it has arbitrarily large integers, because Common
Lisp does.

\paragraph{The compiler to C.}\index{dylanc}
The second is a compiler, \cmd{bin/dylanc}, that translates a complete
Dylan program into C and then runs a C compiler to produce a native
executable. Its source is in \file{compiler/}, and its run-time support is
in \file{runtime/}: a small runtime written in C and the Dylan library,
itself written in Dylan. The compiled programs have proper tail calls,
generic dispatch compiled to decision trees, type inference, inlining,
unboxed floating-point arithmetic, a generational garbage collector,
and real operating-system threads. The system also adds two things the
1992 language does not have: macros, in both implementations, and
threads, in the compiler.

The price of this speed is that the compiler treats a program as
\emph{static}. Everything the program defines must be visible when it is
compiled, and the program cannot change its classes or generic functions
while it runs, cannot evaluate or load new code, and has no bignums or
complex numbers: its integers are 63 bits wide, and overflow is an error.
Chapter~\ref{ch:dylanc} and Appendix~\ref{app:differences} describe these
restrictions in full.

\paragraph{Which to use.} The two are meant to be used together.

\begin{itemize}
\item Use the \textbf{interactive implementation} to learn the language, to
  try out ideas, to explore the library, and to develop code a piece at a
  time. You type a form and see its value at once. It is also the right
  choice for programs that really do need to redefine classes or add
  methods while they run, or that need integers larger than 63 bits.
\item Use the \textbf{compiler} when you want a standalone program: a fast
  executable that starts instantly, needs no Lisp to run, and can use
  threads. A program developed interactively can normally be compiled
  unchanged, provided it stays within the static-program restrictions.
\end{itemize}

Throughout this book, examples work in both implementations unless the
text says otherwise. Where the two differ, for instance in how they print
an object that has no readable printed form, the text points it out.

\section{A short history}\label{sec:intro-history}

\index{Thomas}
In 1992, shortly after the DIRM appeared, Matt Birkholz, Jim Miller and
Ron Weiss of Digital Equipment Corporation's Cambridge Research Laboratory
wrote \term{Thomas}, an implementation of the language in Scheme. They
wrote it in four weeks, and their stated goals were to learn the language
by implementing it from the book, to help others learn from readable
source, and to run small programs; performance was explicitly not a goal.
They named it Thomas rather than Dylan because they had not received
permission to use the trademark.

This project began as a direct translation of Thomas from Scheme into
Common Lisp. Since then it has moved a long way from that starting point.
The Common Lisp implementation was checked against the DIRM chapter by
chapter and corrected wherever Thomas and the book disagreed (for example,
\code{\#f} and the empty list are now distinct objects, as the book
requires), and it gained modules, keyword checking, sealing and slot
options that Thomas lacked. It was then rebuilt for speed: generic
dispatch goes through caches, known calls are made directly, and common
arithmetic, element access, iteration and slot access are compiled in line.

The compiler to C is new. It was written for this project and shares only
the reader with the Common Lisp implementation (which also runs its macros
while compiling). Thomas's own sources are no longer part of the tree;
each file in \file{src/} still records which Thomas file it grew from, and
Thomas's original 1992 README and release notes are kept as
\file{README.thomas} and \file{RELEASE}.

\section{How this book is organized}\label{sec:intro-organization}

The book has three parts and three appendices.

\paragraph{Part I, Getting Started,} is for reading in order.
This chapter is the first. Chapter~\ref{ch:start} explains what you need
to install, how to start the interactive implementation, and how to
compile and run a first program with \cmd{dylanc}.
Chapter~\ref{ch:tutorial} is a hands-on tour of the language, from
arithmetic to classes, generic functions and conditions, ending with a
small complete program built step by step and then compiled.

\paragraph{Part II, The Language,} is a reference to the language, in the
order a reader usually needs it: syntax (Chapter~\ref{ch:syntax}), control
flow (Chapter~\ref{ch:control}), functions (Chapter~\ref{ch:functions}),
classes (Chapter~\ref{ch:classes}), numbers (Chapter~\ref{ch:numbers}),
collections (Chapter~\ref{ch:collections}), functional operations
(Chapter~\ref{ch:functional}), conditions (Chapter~\ref{ch:conditions}),
modules (Chapter~\ref{ch:modules}), and printing and output
(Chapter~\ref{ch:printing}). These chapters follow the DIRM and note,
section by section, where either implementation differs from it.

\paragraph{Part III, The Implementations,} describes each implementation
as a tool: the Common Lisp implementation (Chapter~\ref{ch:lisp}), the
\cmd{dylanc} compiler and its options (Chapter~\ref{ch:dylanc}), the
extensions---macros and threads
(Chapter~\ref{ch:extensions}), how the compiler works inside
(Chapter~\ref{ch:internals}), and how to write fast programs and measure
them (Chapter~\ref{ch:performance}).

\paragraph{The appendices} collect the differences between this system and
the 1992 manual (Appendix~\ref{app:differences}), a summary of the library
(Appendix~\ref{app:library}), and how to run the test suites and
benchmarks (Appendix~\ref{app:testing}).

\section{Conventions}\label{sec:intro-conventions}

\paragraph{Code.} Dylan code, file names, commands and other text you type
or the computer prints are set in a typewriter face, as in
\code{(+ 1 2)} or \file{src/load-dylan.lisp}. In a syntax description,
italic words such as \meta{name} or \meta{body} stand for something you
supply.

\paragraph{Examples and their results.} Most examples are written as
Dylan source. When an example shows what an expression evaluates to, the
value follows the expression in a Dylan comment that begins with
\code{; =>}:

\begin{dylan}
(+ 1 2)          ; => 3
(list 1 'a "b")  ; => (1 a "b")
\end{dylan}

The comment is not part of the value; you can type the whole line,
comment and all, and Dylan will ignore everything from the semicolon to
the end of the line. When an expression returns more than one value, the
values are separated by commas:

\begin{dylan}
(truncate/ 17 5)  ; => 3, 2
\end{dylan}

When an example writes output, the text it writes is shown in a comment
beginning \code{; prints}:

\begin{dylan}
(display "Hello")  ; prints: Hello
\end{dylan}

\paragraph{Terminal sessions.} A session at the shell shows a \code{\$}
prompt before each command you type, followed by the command's output:

\begin{shell}
$ bin/dylanc hello.dyl
$ ./hello
Hello, world!
\end{shell}

\paragraph{Interactive sessions.} A session at the interactive
implementation's read-eval-print loop shows its \code{?} prompt before each
form you type, and the \code{;Value:} line it prints in reply:

\begin{repl}
? (* 6 7)

;Value: 42
\end{repl}

\paragraph{References to the 1992 manual.} A citation such as
\dylanbook{chapter 12} refers to a chapter of the 1992 Dylan manual. An
annotated copy of it, in HTML, is in the repository under
\file{manual/Dylan-1992/}; open \file{START-HERE.html} there. When this
book and the DIRM seem to disagree about the language, the DIRM is the
standard, and any deliberate difference is listed in
Appendix~\ref{app:differences}.

\paragraph{Notes, tips and warnings.} Three kinds of aside appear in the
text. A \textbf{Note} adds a detail or a cross-reference that you can skip
on a first reading. A \textbf{Tip} suggests an easier or better way to do
something. An \textbf{Important} box warns about a mistake that is easy to
make or a difference between the implementations that can surprise you.

\section{License}\label{sec:intro-license}

\index{license}
This software is \emph{source-available}, not open source. The original
Thomas system is Copyright 1992 Digital Equipment Corporation, and the
Common Lisp implementation, which began as a translation of it, is a
derivative work governed by Digital's 1992 notice. That notice allows you
to use, copy, modify and redistribute the software at no charge to
recipients who accept the same terms, but it forbids selling it or
transferring it for compensation. Because of that restriction, the terms
meet neither the Open Source Definition nor the Free Software Definition
and are not compatible with the GPL, MIT, BSD or Apache licenses.

The notice must appear in all copies, derivative works and modified
versions. It is reproduced in full in the file \file{LICENSE} at the top of
the repository, together with an explanation of what it permits; read that
file before redistributing any part of this system.

The home of the project is \url{https://github.com/blakemcbride/Dylan}.
